\subsection{Combustion Modelling}

\subsubsection{Give a short description of the three main approaches to the computation of turbulent flames and specify their advantages and drawbacks.}

{\color{gray}\footnotesize [Claude context: check Lecture 10 (overview of combustion models -- PDF methods, flamelet models, mean reaction rates); Warnatz Ch.13; combustion\_modeling.pdf (supplementary, verify against Warnatz)]}

\paragraph{Answer}  same as \ref{Q:3aproax}



\subsubsection{Explain the main closure problems that arise when averaging the transport equations for an exothermic reacting flow.}

{\color{gray}\footnotesize [Claude context: check Lecture 10; Warnatz Ch.13 (Turbulent Reacting Flows); combustion\_modeling.pdf (supplementary, verify against Warnatz)]}

\paragraph{Answer} same as \ref{Q:clouseure}




\subsubsection{Density-weighted averaging, also called Favre averaging, is an alternative way of averaging leading to a specific form of the averaged transport equations. What is Favre averaging? What is this specific form?}

{\color{gray}\footnotesize [Claude context: check Lecture 10; Warnatz Ch.13 (Favre/density-weighted averaging); combustion\_modeling.pdf (supplementary, verify against Warnatz)]}

\paragraph{Answer} same as \ref{Q:density}




\subsubsection{Explain what is meant by Large Eddy Simulation of turbulent flow. Also discuss the need for subgrid models of combustion processes in the context of LES.}

{\color{gray}\footnotesize [Claude context: check Lecture 10; Warnatz Ch.13 (LES, subgrid closures); combustion\_modeling.pdf (supplementary, verify against Warnatz)]}

\paragraph{Answer} same as \ref{Q:LES}




\subsubsection{Define mixture fraction. Which conditions have to be satisfied in order that mixture fraction becomes a useful variable? Explain the idea that, using mixture fraction, a non-premixed flame problem can be split into a mixing problem and a chemical problem.}

{\color{gray}\footnotesize [Claude context: check Lecture 9/10; Maas notes Ch.13 / Warnatz Ch.14 (Turbulent Non-Premixed Flames); combustion\_modeling.pdf (supplementary, verify against Warnatz/Maas)]}

\paragraph{Answer} same as \ref{Q:mixfrac}




\subsubsection{Explain the structure of the strained laminar flamelet model for non-premixed combustion.}

{\color{gray}\footnotesize [Claude context: check Lecture 10; Warnatz Ch.14 (flamelet model, strained flames); Dynamics of Stretched Flames (Law, 1988, supplementary on flame stretch, verify against Warnatz); combustion\_modeling.pdf (supplementary, verify against Warnatz)]}

\paragraph{Answer} same as \ref{Q:stedylaminatflameletmodel}




\subsubsection{Explain how the influence of turbulent fluctuations of mixture fraction can be taken into account by making use of an assumed shape of its probability density function.}

{\color{gray}\footnotesize [Claude context: check Lecture 10; Warnatz Ch.14 (presumed PDF approach); combustion\_modeling.pdf (supplementary, verify against Warnatz)]}

\paragraph{Answer}

For turbulent floe equations have to be writen in termos of a fable average ( density average). Nevertheless the influence of turbulence still has to be taken into acount. In onder to determine teh aveage influence of turbulence a prescribed PDF is used. This pdf allows the calcualtion of an average though an integral, allowing the determination of the average values. 


similar to \ref{Q:AssumedPDF}




\subsubsection{Explain why one is interested in calculating the variance of mixture fraction. How does the scalar dissipation rate influence the variance of mixture fraction?}

{\color{gray}\footnotesize [Claude context: check Lecture 10; Warnatz Ch.14 (mixture fraction variance, scalar dissipation rate); combustion\_modeling.pdf (supplementary, verify against Warnatz)]}

\paragraph{Answer} similar to \ref{Q:Variance of mixure fraction}

In order to calcuate the average quantity of a property, including the infkence of turbulent flow and variations a PDF is needed. Such PDF is defined with an asusmed shape, a mean and a varance. Variance is tehrefore inportant for the defenition of the PDF function. 

Scalar dissipation rate quantifies how ,uch turbulence is influencin the mxure fractin gradient in dispersion. X increases dissipation or te variace of teh assumed PDF. Scalar dissipation rate plays a dissipation role in teh equation of variance. 




\subsubsection{Explain how model development and experiments come together in the process of model validation. Provide some conditions for good model validation.}

{\color{gray}\footnotesize [Claude context: check Lecture 10; Lecture 5/6 (experimental methods, canonical flames); combustion\_modeling.pdf (validation against DNS/experiments, supplementary, verify against Warnatz where possible)]}

\paragraph{Answer}  many models need empirica constants

Experimetns can detemine the dependent variables

Validation, for same conditions




\subsubsection{List some key characteristics of the Sandia Flame series (D-E-F).}

{\color{gray}\footnotesize [Claude context: check Lecture 10; combustion\_modeling.pdf (TNF workshop benchmark flames, supplementary -- explicitly flag as unverified against infallible sources if not covered in Warnatz/Maas)]}

\paragraph{Answer}

They are increasing in speed






\subsubsection{Discuss the presence and importance of stretch in the counterflow flame configuration.}

{\color{gray}\footnotesize [Claude context: check Lecture 9; Dynamics of Stretched Flames (Law, 1988, supplementary -- explicitly flag as unverified against infallible sources if not covered in Warnatz/Maas); Warnatz Ch.7/Ch.13 for flame structure]}

\paragraph{Answer} The most inportant is tangencial tretch. The tangential stretch caries the radicals outwarsds, possibly extingishing the falem front. 




\subsubsection{How can a laminar premixed flame be described in an intrinsic way using a so-called progress variable? How can a turbulent premixed flame be described using a progress variable?}

{\color{gray}\footnotesize [Claude context: check Lecture 9/10; Warnatz Ch.15 (progress variable, premixed turbulent combustion); combustion\_modeling.pdf (supplementary, verify against Warnatz)]}

\paragraph{Answer} same as \ref{Q:premix progress var}






\subsubsection{Provide an outline of the FGM approach to turbulent non-premixed combustion.}

{\color{gray}\footnotesize [Claude context: check Lecture 10; combustion\_modeling.pdf (FGM/flamelet-generated manifolds, supplementary -- explicitly flag as unverified against infallible sources if not covered in Warnatz/Maas)]}

\paragraph{Answer} Not sure what FGM is

\todo{for Claude: still write the awnser to this one}




\subsubsection{In which sense is the Transported PDF method to turbulent combustion completely different from flamelet models?}

{\color{gray}\footnotesize [Claude context: check Lecture 10; Warnatz Ch.13 (PDF methods overview); combustion\_modeling.pdf (supplementary, verify against Warnatz)]}

\paragraph{Answer} They are 2 difernt models.
Flamelet should be used when flame is much finer than any turbulent scale. Considering a thin reaction flame and negleting trasport on the tangential direction simplifies teh simulation. In this simulation, near the flame cordinates with Z and their normal directions are used. \todo{For calude: is this only for non premixed? Is PDF only for premixed? }

PDF is ment for when turbulenc and chemestry canot be decupled. Turbulent influences mixing. Turbulence requires the usag of density average quantities, especially for areenyus and the chemical sorce term canot be evaluated simply by a mean quantity. To solve this teh average influnce of teh turbulent fluctuations is calcualted from a PDF




\subsubsection{In which cases are flamelet models recommended, and in which cases transported PDF models?}

{\color{gray}\footnotesize [Claude context: check Lecture 10; Warnatz Ch.13/Ch.14; combustion\_modeling.pdf (supplementary, verify against Warnatz)]}

\paragraph{Answer} based on previus awnser and the peter diagram

Flamelets when sl> u rms or ka < 1 - t chem < t k . question  \ref{Q: models peters}





